% !TEX program = xelatex
% ============================================================
%  Arabic RACI Matrix Template — Nibras Center
%  مشروع بحثي: مصفوفة المسؤوليات
%  R = Responsible (منفّذ) | A = Accountable (مسؤول)
%  C = Consulted (مُستشار) | I = Informed (مُطّلع)
%  Compile with: xelatex main.tex (twice)
% ============================================================

\documentclass[11pt,a4paper]{article}

\usepackage[a4paper,landscape,margin=2cm,top=2.5cm,bottom=2.5cm,headheight=15pt]{geometry}
\usepackage{amsmath}
\usepackage{array}
\usepackage{booktabs}
\usepackage{tabularx}
\usepackage{multirow}
\usepackage{enumitem}
\usepackage{float}
\usepackage{titlesec}
\usepackage{fancyhdr}
\usepackage{setspace}
\usepackage[table]{xcolor}

\usepackage{bidi}
\usepackage[no-math]{fontspec}
\usepackage[quiet]{polyglossia}

\setmainfont[Scale=1]{Amiri-Regular.ttf}
\newfontfamily\arabicfont[Script=Arabic,Scale=0.95]{Amiri-Regular.ttf}
\newfontfamily\englishfont[Scale=0.95]{TeX Gyre Termes}

\setdefaultlanguage[calendar=gregorian,hijricorrection=1]{arabic}
\setotherlanguage[variant=british]{english}

\usepackage[breaklinks,hidelinks]{hyperref}
\hypersetup{pdftitle={مصفوفة RACI},pdfauthor={الباحث الرئيسي}}

\titleformat{\section}{\Large\bfseries}{\thesection}{0.7em}{}
\titlespacing*{\section}{0pt}{1.4ex plus 0.4ex}{0.9ex plus 0.3ex}

\pagestyle{fancy}
\fancyhf{}
\fancyhead[R]{\small\itshape مصفوفة المسؤوليات (RACI)}
\fancyhead[L]{\small اسم المشروع}
\fancyfoot[C]{\small\thepage}
\renewcommand{\headrulewidth}{0.3pt}

\newcolumntype{L}[1]{>{\raggedright\arraybackslash}p{#1}}
\newcolumntype{C}[1]{>{\centering\arraybackslash}p{#1}}

\definecolor{raciR}{HTML}{C0392B}
\definecolor{raciA}{HTML}{1A3C6B}
\definecolor{raciC}{HTML}{D4A017}
\definecolor{raciI}{HTML}{27AE60}

\setstretch{1.3}

\begin{document}

\begin{center}
{\LARGE\bfseries مصفوفة المسؤوليات (RACI)}\\[0.5em]
{\Large عنوان المشروع البحثي}\\[1em]
\end{center}

% ============================================================
% LEGEND
% ============================================================
\begin{center}
\renewcommand{\arraystretch}{1.5}
\begin{tabular}{|c|l|}
\hline
\textbf{\color{raciR}R} & \textbf{منفّذ} (\LR{Responsible}) — ينفذ المهمة فعلياً \\
\hline
\textbf{\color{raciA}A} & \textbf{مسؤول} (\LR{Accountable}) — يتحمل المسؤولية النهائية \\
\hline
\textbf{\color{raciC}C} & \textbf{مُستشار} (\LR{Consulted}) — يُستشار قبل القرار \\
\hline
\textbf{\color{raciI}I} & \textbf{مُطّلع} (\LR{Informed}) — يُطّلع بعد القرار \\
\hline
\end{tabular}
\end{center}

\vspace{1em}

% ============================================================
% ROLES
% ============================================================
\section{الأدوار}
\label{sec:roles}

\begin{table}[H]
\centering
\begin{tabularx}{\textwidth}{C{1.5cm} L{5cm} X}
\toprule
\textbf{الرمز} & \textbf{الدور} & \textbf{الاسم} \\
\midrule
PI & الباحث الرئيسي (\LR{Principal Investigator}) & --- \\
Co1 & باحث مشارك أول & --- \\
Co2 & باحث مشارك ثانٍ & --- \\
Tech & فني مختبر & --- \\
PM & مدير المشروع & --- \\
Fund & ممثل الجهة الممولة & --- \\
\bottomrule
\end{tabularx}
\end{table}

% ============================================================
% RACI MATRIX
% ============================================================
\section{مصفوفة RACI}
\label{sec:raci_matrix}

\begin{table}[H]
\centering
\small
\renewcommand{\arraystretch}{1.4}
\begin{tabularx}{\textwidth}{L{5cm} C{1.2cm} C{1.2cm} C{1.2cm} C{1.2cm} C{1.2cm} C{1.2cm}}
\toprule
\textbf{المهمة / النشاط} & \textbf{PI} & \textbf{Co1} & \textbf{Co2} & \textbf{Tech} & \textbf{PM} & \textbf{Fund} \\
\midrule
\rowcolor{lightgray!20}
\multicolumn{7}{l}{\textbf{1. المراجعة والتخطيط}} \\
مراجعة الأدبيات & A/R & C & R & --- & I & I \\
تحديد الفجوات البحثية & A/R & C & C & --- & I & --- \\
تصميم المنهجية & A/R & C & C & --- & I & I \\
إعداد ميثاق المشروع & A & C & --- & --- & R & I \\
\midrule
\rowcolor{lightgray!20}
\multicolumn{7}{l}{\textbf{2. الإعداد والتجربة}} \\
تجهيز المعدات & A & C & --- & R & I & --- \\
جمع البيانات & A & R & R & C & I & --- \\
معالجة البيانات & A & R & C & --- & I & --- \\
\midrule
\rowcolor{lightgray!20}
\multicolumn{7}{l}{\textbf{3. التحليل والنمذجة}} \\
التحليل الإحصائي & A/R & C & C & --- & I & --- \\
النمذجة والمحاكاة & A & R & R & --- & I & --- \\
التحقق من النتائج & A/R & C & C & --- & I & I \\
\midrule
\rowcolor{lightgray!20}
\multicolumn{7}{l}{\textbf{4. النشر والتوثيق}} \\
كتابة الورقة البحثية & A/R & C & C & --- & I & --- \\
تقديم الورقة للنشر & A/R & I & I & --- & I & I \\
كتابة التقرير النهائي & A/R & C & C & --- & C & I \\
عرض النتائج & A/R & R & R & C & C & I \\
إغلاق المشروع & A & C & C & --- & R & C \\
\bottomrule
\end{tabularx}
\end{table}

% ============================================================
% NOTES
% ============================================================
\section{ملاحظات}
\label{sec:notes}

\begin{itemize}
    \item لكل مهمة يجب أن يكون هناك \textbf{مسؤول واحد فقط (A)}.
    \item يمكن أن يكون هناك عدة \textbf{منفّذين (R)} للمهمة الواحدة.
    \item \textbf{المُستشار (C)} يُطلب رأيه ثنائياً قبل اتخاذ القرار.
    \item \textbf{المُطّلع (I)} يُبلغ بالقرار بعد اتخاذه.
    \item إذا كان الشخص \textbf{A/R} فهو مسؤول ومنفّذ في آن واحد.
\end{itemize}

\end{document}
